% !TeX encoding = UTF-8
% !TeX program = xelatex
\documentclass[type=bachelor,oneside,openany]{sjtuthesis}

% !TEX root = ./main.tex

\sjtusetup{
  info = {
    zh / title      = {鞅过程可选停止定理的数值模拟},
    en / title      = {Numerical Simulation of the Optional Stopping Theorem for Martingales},
    zh / subject    = {随机过程课程大作业},
    en / subject    = {Stochastic Processes Course Project},
    zh / keywords   = {可选停止定理, 鞅, Monte Carlo 模拟, 最优停时, 停时, Snell 包络},
    en / keywords   = {optional stopping theorem, martingale, Monte Carlo simulation, optimal stopping, stopping time, Snell envelope},
    zh / author     = {杨逸飞},
    en / author     = {Yifei Yang},
    id              = {522111910124},
    zh / department = {化学化工学院},
    en / department = {School of Chemistry and Chemical Engineering},
    zh / major      = {随机过程课程大作业},
    en / major      = {Stochastic Processes Course Project},
    zh / degree     = {课程作业},
    en / degree     = {Course Project},
    date            = {2026-06-01},
  },
  style = {
    keywords-format = hang,
  },
  name = {
    bib = {参考文献},
  },
}

\usepackage[backend=biber,style=gb7714-2015,gbnamefmt=uppercase,maxbibnames=3,minbibnames=3,doi=false]{biblatex}
\renewcommand{\bibfont}{\zihao{5}\setbaselineskip{16bp}}
\setlength{\bibitemsep}{3bp plus 1pt}
\addbibresource{references.bib}

\usepackage[perpage,bottom,hang]{footmisc}
\usepackage{graphicx}
\graphicspath{{figures/}}
\DeclareGraphicsExtensions{.pdf,.eps,.png,.jpg,.jpeg}
\usepackage{booktabs}
\usepackage{longtable}
\usepackage{float}
\usepackage{amsthm}
\usepackage{chngcntr}
\usepackage{indentfirst}
\captionsetup[figure]{justification=centering}

\ctexset{
  section = {
    format = \centering\zihao{2}\bfseries,
    number = \arabic{section},
    name = {第~,~章},
    aftername = \quad,
    beforeskip = 0pt,
    afterskip = 24bp,
  },
  subsection = {
    format = \zihao{3}\bfseries,
    beforeskip = 18bp,
    afterskip = 10bp,
  },
  subsubsection = {
    format = \zihao{-3}\bfseries,
    beforeskip = 12bp,
    afterskip = 6bp,
  },
}

\contentsmargin[2.55em]{0pt}
\titlecontents{section}
  [0pt]
  {\addvspace{6bp}\bfseries}
  {\contentspush{\thecontentslabel\enskip}}
  {}
  {\titlerule*[4bp]{.}\contentspage}

\titlecontents{subsection}
  [2em]
  {}
  {\contentspush{\thecontentslabel\enskip}}
  {}
  {\titlerule*[4bp]{.}\contentspage}

\titlecontents{subsubsection}
  [4em]
  {}
  {\contentspush{\thecontentslabel\enskip}}
  {}
  {\titlerule*[4bp]{.}\contentspage}

\theoremstyle{definition}
\newtheorem{definition}{定义}[section]

\theoremstyle{plain}
\newtheorem{theorem}{定理}[section]
\newtheorem{lemma}[theorem]{引理}
\newtheorem{corollary}[theorem]{推论}

\theoremstyle{remark}

\setcounter{tocdepth}{3}


\newcommand{\chapterinput}[1]{\clearpage\input{#1}}

\newcommand{\makecoursecover}{
  \begin{titlepage}
    \thispagestyle{empty}
    \centering
    \vspace*{2.8cm}
    {\zihao{-1}\bfseries 上海交通大学\par}
    \vspace{1.8cm}
    {\zihao{2}\bfseries 鞅过程可选停止定理的数值模拟\par}
    \vspace{1.0cm}
    {\zihao{3}\bfseries 随机过程课程大作业\par}
    \vfill
    {\zihao{4}
    \begin{tabular}{rl}
      作者： & 杨逸飞 \\
      学号： & 522111910124 \\
      学院： & 化学化工学院 \\
      班级： & 化工2202 \\
    \end{tabular}\par}
    \vfill
    {\zihao{4} 2026 年 6 月\par}
  \end{titlepage}
}

\begin{document}

\makecoursecover

\frontmatter

\begin{abstract}[zh]
可选停止定理（Optional Stopping Theorem, OST）是鞅论中的核心结论之一，它刻画了随机过程中”在停时处停止”与”在固定时刻观察”之间的期望关系。本文采用理论推导与 Monte Carlo 数值实验相结合的方法，系统讨论 OST 在不同模型中的成立条件与失效机制。

本文选取三个具有代表性的随机过程模型作为分析对象：Moran 种群遗传过程、Cram\'{e}r-Lundberg 保险破产过程和赌徒破产过程。这三个模型分别对应 OST 理论谱系中的三种典型情形：Moran 模型展示了 OST 在有界状态空间中直接成立的情形；保险破产模型展示了如何通过指数变换构造辅助鞅来应用 OST（且只能得到不等式上界）；赌徒破产模型则通过双边与单边障碍的对比，展示了当停时无界、过程无界且停时期望发散时 OST 如何失效。

在数值方法上，本文统一构建了 Monte Carlo 模拟框架，对各模型分别执行路径生成、停时记录、期望估计与图表可视化。结果表明：Moran 模型中的固定概率与理论公式 $x_0/N$ 高度吻合；保险破产模型中的模拟结果验证了 Lundberg 上界与指数衰减规律；赌徒破产模型直观呈现了截断停时与自然停时在期望上的本质差异，并通过停时分布的尾部行为揭示了一致可积性缺失的微观机制。

综合来看，本文不仅展示了 OST 在不同随机过程中的统一解释力，也揭示了其适用范围的边界。三个模型在停时分析的难度上从简单到复杂依次递进，构成了从”OST 直接成立”到”OST 条件失效”的完整谱系。本文力求在数学逻辑、实验设计与图表表达三个层面形成较完整的分析链条，为理解鞅、停时及 OST 提供可计算、可验证的案例集合。
\end{abstract}

\tableofcontents
\listoffigures
\listoftables
% !TEX root = ./main.tex

\begin{nomenclature}
\label{chap:symb}

\begin{longtable}{p{0.22\textwidth}p{0.68\textwidth}}
OST & 可选停止定理（Optional Stopping Theorem） \\
$\mathcal{F}_n,\mathcal{F}_t$ & 到时刻 $n$ 或 $t$ 为止的信息流（滤过） \\
$\tau$ & 停时 \\
$X_n, X_t$ & 一般随机过程或鞅过程记号 \\
$S_n$ & 随机游走或赌徒净收益过程 \\
$N$ & Moran 模型中的种群规模 \\
$x_0$ & Moran 模型中 A 等位基因的初始个数 \\
$\theta$ & 保险模型中的安全负荷系数，$\theta = c/(\lambda\mu)-1$ \\
$\lambda$ & 索赔到达强度（泊松过程参数） \\
$\mu$ & 索赔额均值 $\mathbb{E}[Y]$ \\
$c$ & 保费速率 \\
$a, b$ & 赌徒破产模型中甲、乙的初始资金 \\
$p, q$ & 随机游走中每一步赢/输的概率，$q = 1-p$ \\
$P_N$ & Moran 模型中 A 最终固定的概率 \\
$U_t$ & 保险盈余过程 \\
$\psi(u)$ & 初始资本为 $u$ 时的最终破产概率 \\
$R$ & 保险破产模型中的调整系数 \\
$M$ & Monte Carlo 模拟路径数 \\
\end{longtable}

\end{nomenclature}


\mainmatter
\counterwithout{section}{chapter}
\counterwithout{equation}{chapter}
\counterwithout{figure}{chapter}
\counterwithout{table}{chapter}
\counterwithin{equation}{section}
\counterwithin{figure}{section}
\counterwithin{table}{section}
\renewcommand{\thesection}{\arabic{section}}
\renewcommand{\thesubsection}{\thesection.\arabic{subsection}}
\renewcommand{\thesubsubsection}{\thesubsection.\arabic{subsubsection}}
\renewcommand{\sectionmark}[1]{\markboth{第\,\thesection\,章\ #1}{第\,\thesection\,章\ #1}}

\chapterinput{sections/ch01_intro}
\chapterinput{sections/ch02_moran}
\chapterinput{sections/ch03_insurance}
\chapterinput{sections/ch04_ruin}
\chapterinput{sections/ch07_conclusion}

\backmatter
\printbibliography[heading=bibintoc]

\end{document}
